% from ega3-1-fr.tex
\oldpage[0\textsubscript{III}]{29}
Supposons que dans $\mathcal C$ les limites inductives filtrantes soient
exactes. Alors, si la filtration $(F^p(K^\bullet))$ de $K^\bullet$ est
\emph{exhaustive}, il en est de même de la filtration
$(F^p(\mathrm H^n(K^\bullet)))$ pour tout $n$, car on a par hypothèse
$K^\bullet=\varinjlim_{p\to-\infty}F^p(K^\bullet)$ et l'hypothèse sur
$\mathcal C$ entraîne que la cohomologie commute aux limites inductives. En
outre, on a, pour la même raison,
$B_\infty(E_2^{pq}(K^\bullet))=\sup_k B_k(E_2^{pq}(K^\bullet))$. On dit que
la filtration $(F^p(K^\bullet))$ de $K^\bullet$ est \emph{régulière} si pour
tout $n$ il existe un entier $u(n)$ tel que
$\mathrm H^n(F^p(K^\bullet))=0$ pour $p>u(n)$. Il en est ainsi en particulier
lorsque la filtration de $K^\bullet$ est \emph{discrète}. Lorsque la
filtration de $K^\bullet$ est régulière et exhaustive, et que les limites
inductives filtrantes dans $\mathcal C$ sont exactes, on montre (M, XV, 4)
que la suite spectrale $E(K^\bullet)$ est \emph{régulière}.
\end{env}

\subsection{Les suites spectrales d'un bicomplexe}
\label{subsection:0.11.3-fr}

\begin{env}[11.3.1]
\label{0.11.3.1-fr}
En ce qui concerne les conventions relatives aux bicomplexes, nous suivons
celles de (T, 2.4) plutôt que celles de (M), les deux dérivations $d'$, $d''$
(de degré $+1$) d'un tel bicomplexe
$K^{\bullet\bullet}=(K^{ij})$ étant donc supposées \emph{permutables}.
Supposons vérifiées l'une des deux conditions suivantes~:
\begin{enumerate}
  \item[$1^\circ$] Les sommes directes infinies existent dans $\mathcal C$.
  \item[$2^\circ$] Pour tout $n\in\mathbb Z$, il n'y a qu'un nombre
    \emph{fini} de couples $(p,q)$ tels que $p+q=n$ et $K^{pq}\ne0$.
\end{enumerate}
Alors, le bicomplexe $K^{\bullet\bullet}$ définit un complexe (simple)
$(K^{\prime n})_{n\in\mathbb Z}$, avec
$K^{\prime n}=\sum_{i+j=n}K^{ij}$, l'opérateur de dérivation $d$ (de degré
$+1$) de ce complexe étant donné par
$dx=d'x+(-1)^i d''x$ pour $x\in K^{ij}$. Lorsque nous parlerons par la suite
du complexe (simple) \emph{défini par un bicomplexe $K^{\bullet\bullet}$}, il
sera toujours sous-entendu que l'une des conditions précédentes est satisfaite.
On adoptera des conventions analogues pour les multicomplexes.

On désigne par $K^{i,\bullet}$ (resp. $K^{\bullet,j}$) le complexe simple
$(K^{ij})_{j\in\mathbb Z}$ (resp. $(K^{ij})_{i\in\mathbb Z}$), par
$Z_{\mathrm{II}}^p(K^{i,\bullet})$, $B_{\mathrm{II}}^p(K^{i,\bullet})$,
$\mathrm H_{\mathrm{II}}^p(K^{i,\bullet})$ (resp.
$Z_{\mathrm I}^p(K^{\bullet,j})$, $B_{\mathrm I}^p(K^{\bullet,j})$,
$\mathrm H_{\mathrm I}^p(K^{\bullet,j})$) ses $p$-èmes objets des cocycles,
des cobords et de cohomologie respectivement~; la dérivation
$d':K^{i,\bullet}\to K^{i+1,\bullet}$ est un morphisme de complexes, qui
donne donc un opérateur dans les cocycles, les cobords et la cohomologie,
\[
\begin{aligned}
  d'&:Z_{\mathrm{II}}^p(K^{i,\bullet})
      \longrightarrow Z_{\mathrm{II}}^p(K^{i+1,\bullet}),\\
  d'&:B_{\mathrm{II}}^p(K^{i,\bullet})
      \longrightarrow B_{\mathrm{II}}^p(K^{i+1,\bullet}),\\
  d'&:\mathrm H_{\mathrm{II}}^p(K^{i,\bullet})
      \longrightarrow \mathrm H_{\mathrm{II}}^p(K^{i+1,\bullet}),
\end{aligned}
\]
et il est clair que pour ces opérateurs,
$(Z_{\mathrm{II}}^p(K^{i,\bullet}))_{i\in\mathbb Z}$,
$(B_{\mathrm{II}}^p(K^{i,\bullet}))_{i\in\mathbb Z}$ et
$(\mathrm H_{\mathrm{II}}^p(K^{i,\bullet}))_{i\in\mathbb Z}$ sont des
complexes~; nous désignerons le complexe
$(\mathrm H_{\mathrm{II}}^p(K^{i,\bullet}))_{i\in\mathbb Z}$ par
$\mathrm H_{\mathrm{II}}^p(K^{\bullet\bullet})$, ses $q$-èmes objets de
cocycles, de cobords et de cohomologie par
$Z_{\mathrm I}^q(\mathrm H_{\mathrm{II}}^p(K^{\bullet\bullet}))$,
$B_{\mathrm I}^q(\mathrm H_{\mathrm{II}}^p(K^{\bullet\bullet}))$ et
$\mathrm H_{\mathrm I}^q(\mathrm H_{\mathrm{II}}^p(K^{\bullet\bullet}))$.
On définit de même les complexes
$\mathrm H_{\mathrm I}^p(K^{\bullet\bullet})$ et leurs objets de cohomologie
$\mathrm H_{\mathrm{II}}^q(\mathrm H_{\mathrm I}^p(K^{\bullet\bullet}))$.
Rappelons d'autre part que $\mathrm H^n(K^{\bullet\bullet})$ désigne le
$n$-ème objet de cohomologie du complexe (simple) défini par
$K^{\bullet\bullet}$.
\end{env}

\begin{env}[11.3.2]
\label{0.11.3.2-fr}
Sur le complexe défini par un bicomplexe $K^{\bullet\bullet}$, on peut
considérer deux filtrations canoniques
$(F_{\mathrm I}^p(K^{\bullet\bullet}))$ et
$(F_{\mathrm{II}}^p(K^{\bullet\bullet}))$ données par
\[
\label{0.11.3.2.1-fr}
  F_{\mathrm I}^p(K^{\bullet\bullet})
  =\left(\sum_{\substack{i+j=n\\i\geq p}}K^{ij}\right)_{n\in\mathbb Z},
  \qquad
  F_{\mathrm{II}}^p(K^{\bullet\bullet})
  =\left(\sum_{\substack{i+j=n\\j\geq p}}K^{ij}\right)_{n\in\mathbb Z}.
  \tag{11.3.2.1}
\]
